\documentclass[a4paper,10pt]{article}

\usepackage{../preamble-paper}
\addbibresource{../ref.bib}
\AtEveryBibitem{\clearfield{note}}

\fancyhead[R]{\thepage}
\fancyhead[L]{\textbf{交易研究}}

\begin{document}

\title{\textbf{交易研究}\\ \large }
\author{Lao Chen}
\date{\today}
\maketitle
\tableofcontents

\section{随机过程}

从价格序列推断其背后的随机过程模型，是金融计量经济学的核心议题。
近年来的研究呈现出两条清晰的主线：
一条是在传统随机过程模型框架内持续深耕，
另一条是利用机器学习方法开辟新的路径。

\cite{MaungSwanson2025Forecasting} 
对当前方法进行了较全面的概述，
系统地回顾了从简单的随机游走模型，
到离散时间的 GARCH 族模型，
再到连续时间的随机波动率模型，
并讨论了大数据和机器学习方法在其中的应用，
直接回应了从数据中建模和预测金融时间序列的各种工具。

\cite{Chakraborti2007FinTSOverview} 
则清晰地梳理了金融时间序列的几个关键“典型事实”（stylized facts），
如厚尾（fat tails）、
波动率聚集（volatility clustering）和收益率的自相关缺失，
并系统介绍了用于描述这些特征的随机模型
（如随机游走、GARCH、跳跃扩散模型），
是理解为何需要复杂模型的基础。

在机器学习与高维方法方面，
\cite{Yang2025MLFinTS} 
专门聚焦于机器学习方法，
如神经网络、随机森林、LASSO 回归等在金融时间序列预测中的应用，
讨论了如何利用这些方法处理金融数据的非线性、动态性和混沌特性，
并比较了不同模型的性能。

\cite{Masini2023MLAdvances} 
则深入探讨了监督式机器学习的最新进展，
特别是惩罚回归、集成方法以及浅层和深度神经网络（包括前馈和循环神经网络），
不仅介绍了方法，还讨论了在金融和经济学中的应用和预测能力检验。

\cite{EmmertStreib2026TSFMethods} 
的覆盖面更广，从传统统计模型一直延伸到深度学习架构，
如 RNN、CNN、Transformer、图神经网络，甚至包括 Neural ODEs 等较新的方法，
为模型选择提供了一个全景式的分类和比较框架。

在特定模型的深度研究上，
\cite{Virbickaite2015BayesianGARCH}
专注于 GARCH 模型的贝叶斯推断方法，
详细比较了各种参数和非参数贝叶斯方法，
特别是讨论了对厚尾分布的灵活处理。

\cite{Sims2010ContinuousTime} 
则从理论层面讨论了金融中连续时间随机过程
（如几何布朗运动、Cox-Ingersoll-Ross 模型等）
的设定与估计问题，
为理论模型如何与实际数据结合提供了视角。

总体来看，当前该领域的研究趋势可以概括为：
传统模型（从随机游走到 GARCH、随机波动率模型及其各种扩展，
如 EGARCH、GJR）仍然是为波动率、风险等提供基准预测的核心工具；
而处理非线性、高维数据的需求，
使得深度学习和集成学习方法在预测准确性方面展现出巨大潜力，
成为主流研究范式之一；
同时，高维函数型时间序列分析、贝叶斯非参数方法等新方向，
也为更精细地利用高频、高维数据提供了新的建模思路。

\section{收益率}

在金融时间序列分析中，原始价格序列通常是非平稳的，直接建模会导致伪回归等问题。因此，通常将价格序列转换为收益率序列，使其具有平稳性，便于后续的统计建模和预测。
设 $P_t$ 为资产在时刻 $t$ 的价格（通常使用调整收盘价），$P_{t-1}$ 为前一时刻的价格。
简单收益率和对数收益率定义为：

\begin{equation}
    R_t = \frac{P_t - P_{t-1}}{P_{t-1}} = \frac{P_t}{P_{t-1}} - 1
    \qquad
    r_t = \ln\left(\frac{P_t}{P_{t-1}}\right) = \ln P_t - \ln P_{t-1}
    \label{eq:log_return}
\end{equation}

计算得到的收益率序列通常满足平稳性。
可通过 ADF 检验（Augmented Dickey-Fuller Test）验证：

\begin{equation}
    \Delta r_t = \alpha + \beta t + \gamma r_{t-1} + \sum_{i=1}^{p} \delta_i \Delta r_{t-i} + \varepsilon_t
\end{equation}

若拒绝单位根原假设 $(\gamma = 0)$，则认为序列平稳。
给定初始价格 $P_0$ 和对数收益率序列 $\{r_1, \dots, r_T\}$，可还原价格：
从价格序列转换为收益率序列是金融建模的基础操作。
在实证研究中，推荐使用对数收益率，
因其具有良好的统计性质（时间可加性、近似对称性、近似正态性），
是GARCH、SV等波动率模型的输入标准。

GARCH 和 SV 模型本身不直接基于Brown运动，
它们使用的是离散时间的独立同分布冲击。
但当采样频率提高时，
两者的连续时间极限都会收敛到由Brown运动驱动的随机微分方程。
因此，可以说它们是对Brown运动驱动过程的离散时间近似。
这里是指对数价格空间上的Brown运动。

\subsection{GARCH模型}

设收益率序列 $\{y_t\}$，其均值方程为：
\begin{equation}
    y_t = \mu + \varepsilon_t, \qquad \varepsilon_t = \sigma_t z_t, \qquad z_t \sim \text{i.i.d.}(0,1).
\end{equation}
其中 $\mu$ 为条件均值（可包含 ARMA 结构），$\{\varepsilon_t\}$ 为残差序列。

条件方差 $\sigma_t^2$ 服从如下动态过程：
\begin{equation}
    \sigma_t^2 = \omega + \sum_{i=1}^{q} \alpha_i \varepsilon_{t-i}^2 + \sum_{j=1}^{p} \beta_j \sigma_{t-j}^2, \label{eq:garch_variance}
\end{equation}
其中参数满足：
\begin{equation}
    \omega > 0, \qquad \alpha_i \ge 0 \ (i=1,\dots,q), \qquad \beta_j \ge 0 \ (j=1,\dots,p).
\end{equation}

为保证平稳性，要求：
\begin{equation}
    \sum_{i=1}^{q} \alpha_i + \sum_{j=1}^{p} \beta_j < 1.
\end{equation}
此时无条件方差为：
\begin{equation}
    \operatorname{Var}(y_t) = \frac{\omega}{1 - \sum_{i=1}^{q} \alpha_i - \sum_{j=1}^{p} \beta_j}.
\end{equation}

\subsection{随机波动率}

综合当前实证研究来看，
对于一般的金融价格时间序列，
随机波动率（SV）模型在白箱模型中更接近真实市场 
\cite{PangZhao2025GARCHSV}。
在一系列严格的实证比较中，
SV模型在刻画波动率方面比GARCH模型展现出了更明显的优势。

GARCH和SV模型是白箱框架下最主流的两类选择，
但二者的工作方式有本质区别：
GARCH模型的波动率是过去收益率和过去波动率的确定性函数；
而SV模型的波动率本身有自己的随机过程，
引入了额外的新息，
因此自由度更高，能更灵活地捕捉波动的突变。

实证数据明确支持了SV模型的优势。
在预测精度方面，对标普500指数的一项研究发现，
ARSV(1)模型（一种SV模型）在样本内拟合和样本外预测上表现都优于GARCH(1,1)模型 \cite{PangZhao2025GARCHSV}；在期权定价方面，ARSV模型在拟合和预测看跌期权价格时表现也显著优于GARCH模型 \cite{PangZhao2025GARCHSV}；
在厚尾特征方面，资产收益率普遍存在“尖峰厚尾”特征，能够刻画厚尾特征的SV模型（SV-T）预测能力最好。

金融市场的价格波动本质上是一个受无数新信息冲击的复杂系统，
SV模型内置的随机性更贴合这种现实：
首先，SV模型不把波动率看作过去数据的死结果，
而是认为它本身就是一个受独立随机因素驱动的活过程，
更能解释波动率无法用历史数据完全解释的剧烈变动；
其次，资产收益率的平方或绝对值序列的自相关函数（ACF）通常衰减很慢，
呈现出长记忆性，
SV模型中的某些设定（如长记忆随机波动率模型）能更好地捕捉这一现象；
最后，大量学术研究，包括针对标普500指数和汇率市场的实证分析，
都倾向于支持SV模型在拟合和预测上的优越性 
\cite{PangZhao2025GARCHSV}。

因此，对于“哪个白箱模型更接近真实市场”这一问题，
目前答案最倾向于随机波动率（SV）模型。
虽然GARCH模型因其计算简便仍是业界广泛使用的基准，
但若追求对市场真实波动更精准的刻画以及更可靠的预测能力，
SV模型具有更强的理论和实证依据。

当然，选择模型也需考虑实际约束：
SV模型的参数估计更复杂，
通常需要使用MCMC等模拟方法，
在开发轻量级实时交易系统时可能面临计算成本挑战；
而GARCH模型依然是许多实务场景的首选，
但对于学术研究或构建更精细的风险管理模型，
SV模型值得投入计算资源。

\section{除权}

关于股票除权、拆股（Stock Splits）以及之后的股价表现（包括抗跌性与超额收益），在学术界和行为金融学中确实有非常多深入的研究。

学术界通常不直接使用“抗跌”这个大白话，而是将其放在“名义价格错觉”、“信号理论”和“流动性假说”的框架下去研究。主要的几个核心学术结论如下：

\subsection{1. 名义价格错觉（Nominal Price Illusion）与零碎股冲击}

俄亥俄州立大学的学者 Birru 和 Wang 在 2016 年发表了一篇极具代表性的论文《Nominal Price Illusion》\cite{BirruWang2016}。

\begin{itemize}
    \item \textbf{传统研究发现}：散户投资者对低绝对股价的股票有一种“彩票心理”（Lottery-like stock preference），错误地认为低价股“上涨空间更大，下跌空间有限（更抗跌）”。这是一种典型的认知偏差（Numerosity Bias）。
    \item \textbf{最新研究反转（2024年）}：Borsboom 和 Füllbrunn (2024)\cite{BorsboomFullbrunn2024} 针对新一代互联网券商（Neobrokers）的研究发现，随着\textbf{“零碎股交易（Fractional Shares）”}的普及，这种由名义价格错觉引发的短期“抗跌或暴涨”效应在美股成熟市场已经显著减弱乃至消失。因为资金约束被彻底打破，散户不再因为“买不起一手”而被迫偏好低绝对股价。
    \item \textbf{A股市场的特异性}：然而，由于 \textbf{A股市场依然存在“最小交易单位 100 股”的刚性限制}，这种名义门槛（Affordability Constraint）依然对中国散户具有显著的心理锚定效应，除权低价对散户的吸引力在本土环境依然部分存在。
\end{itemize}

\subsection{2. 信号理论（Signaling Hypothesis）}

著名的金融学者 Ikenberry 等人对拆股/除权进行了长期的跟踪研究（《What do Stock Splits Really Signal?》\cite{Ikenberryetal}）。

\begin{itemize}
    \item \textbf{研究结论}：统计表明，历史上宣布拆股/大比例送转的公司，在随后的 12 个月到 3 年内，往往能跑赢大盘（产生超额收益 CAR）。
    \item \textbf{核心逻辑}：这\textbf{不是因为除权动作本身让他们抗跌}，而是因为“管理层的自选择效应”。管理层只有在极度看好公司未来几年业绩、确定股价不会崩盘的情况下，才敢进行大比例除权。如果管理层预期公司快暴雷了，他们绝不敢除权（否则股价会直接沦为仙股或触发面值退市）。因此，除权实际上是管理层向市场释放的“高分贝信心信号”。
\end{itemize}

\subsection{3. A股市场的本土化研究：“高送转”与贴权陷阱}

针对中国A股的独特生态，国内学术界和券商研究所（如深交所金融创新实验室\cite{SZSE_HighDividend}、各大券商金融工程部）做过专门的量化研究。

\begin{itemize}
    \item \textbf{套利效应的终结}：早期A股存在显著的“高送转预案公告效应”。在除权前，股票往往能获得极高的超额收益。但研究表明，\textbf{在正式除权日（Ex-date）之后，这种超额收益会迅速逆转}。
    \item \textbf{大股东减持与“贴权”}：研究指出，A股相当一部分高比例除权（送转）并不是因为基本面好，而是为了配合大股东解禁减持。在除权后，由于散户盲目涌入接盘，主力资金借机出货，导致股价在除权后表现出极高的\textbf{非抗跌性}（即显著的贴权，表现弱于大盘）。这也是为什么中国证监会后来对“高送转”出台了极其严格的监管指引。
\end{itemize}

\subsection{4. 波动率跃升效应（The Volatility Spike Effect）}

在量化金融学中，评估一只股票是否“抗跌”，核心指标之一是其收益率的波动率（Volatility）。经典文献直接否定了除权能降低风险的假设。

\begin{itemize}
    \item \textbf{实证发现}：Ohlson 和 Penman（1985）\cite{OhlsonPenman1985} 以及 Koski（1995）\cite{Koski1995} 针对美股拆股的大样本量化研究表明，股票在除权（拆股）后，其\textbf{日收益率的波动率（Standard Deviation）和系统性风险（Beta）平均显著上升约 20\% 至 30\%}。
    \item \textbf{定量解释}：这种波动率的上升并不能完全由公司基本面的改变来解释。后续研究表明，这种波动率跃升（Volatility Spike）主要集中在除权后的前两周。由于绝对股价降低，噪音交易者（Noise Traders）和高频散户的频繁买卖引入了更多的市场噪音，反而使股价在面临市场冲击时\textbf{更容易出现剧烈波动，而不是更抗跌}。
\end{itemize}

\subsection{5. 长期非流动性永久恶化（Permanent Illiquidity Deterioration）}

管理层选择除权的重要借口之一是“提高股票流动性并增强股价稳定性”，但最近几年的微观结构前沿研究对这一假说提出了严厉挑战。

\begin{itemize}
    \item \textbf{买卖价差与订单簿受损}：早在 Conroy 等 (1990)\cite{ConroyHarrisBenet1990} 与 Lipson (2001)\cite{Lipson2001} 的基础研究中就发现，除权后限价订单簿（Limit Order Book）的百分比有效价差会扩大，深度变薄 30\%。
    \item \textbf{最新前沿量化（2025年）}：Linton 与 Olive (2025)\cite{LintonOlive2025} 采用最新的动态半参数时间序列模型对近三十年的高频交易数据进行了拆解。实证结果表明，除权对流动性的影响具有严重的“双相特征”：虽然除权在短期内会由于散户订单数量（Order Count）的爆发带来暂时的交易活跃，但**长期来看，股票的阿米胡德非流动性指标（Amihud Illiquidity Trend）呈现出显著且永久性的上升（即流动性恶化）**。
    \item \textbf{击穿效应}：这意味着，长期来看，除权去除了原先高价时机构资金提供的深厚限价单缓冲。当市场发生系统性恐慌抛盘时，变薄的订单簿极易引发“流动性黑洞”与价格踩踏，**导致除权后的股票在熊市或崩盘期间表现出极弱的抗跌性**。
\end{itemize}

\subsection{总结学术界给出的量化视角}

\begin{table}[h]
\centering
\begin{tabular}{|p{3.5cm}|p{6.5cm}|p{4.5cm}|}
\hline
\textbf{研究维度} & \textbf{学术结论} & \textbf{对“抗跌”的实际映射} \\
\hline
\textbf{短期（除权后几天）} & 存在由于散户“名义价格错觉”带来的资金净流入，但在零碎股普及背景下逐步减弱。 & \textbf{微弱抗跌}。在存在 100 股限制的 A 股市场门槛降低，交易量技术性上升。 \\
\hline
\textbf{微观结构（除权后数周）} & 日内波动率与 Beta 显著跃升 20\%--30\%，买卖百分比价差扩大。 & \textbf{不抗跌}。市场噪音增加导致价格日内震荡加剧，容易被极端单边单诱导。 \\
\hline
\textbf{长期微观特征（2025前沿）} & 长期 Amihud 非流动性趋势永久性恶化。机构订单簿深度变薄。 & \textbf{极不抗跌}。缺乏大单缓冲垫，在面临系统性大盘杀跌时极易形成价格击穿。 \\
\hline
\textbf{中长期（除权后数月）} & 表现高度分化。如果是为掩护大股东解禁送转，除权后往往陷入“贴权陷阱”。 & \textbf{更危险}。彻底暴露出缺乏基本面支撑的投机本质。 \\
\hline
\end{tabular}
\end{table}

如果想更直观地通过数据和历史回测了解除权、拆股对股价产生的真实收益影响，推荐观看 \href{https://www.youtube.com/watch?v=MgMTBNq7Bi4}{Do Stock Splits ACTUALLY Boost Returns? What the Numbers Say} 这段视频\cite{YouTube_StockSplits}。视频用大量的历史回测数据客观拆解了“拆股/除权”在公告日和执行日之后的真实统计表现，能帮你从数据上厘清这种“抗跌”或“暴涨”的迷思。

\printbibliography[heading=bibintoc]

\end{document}